\documentclass[aspectratio=169,10pt,spanish]{beamer}

\input{../../../_preambulo_beamer.tex}
\input{../../../_comandos_beamer.tex}

\selectlanguage{spanish}

\title{MA1001B - Modelación Estadística para la Toma de Decisiones}
\subtitle{Lección 46: Supuestos y diagnósticos de ANOVA}
\author{}
\institute{Tecnol\'ogico de Monterrey}
\date{}

\begin{document}

\begin{frame}[plain]
  \vspace{1.2cm}
  {\Large\bfseries MA1001B - Modelación Estadística para la Toma de Decisiones\par}
  \vspace{0.7cm}
  {\large Lección 46: Supuestos y diagnósticos de ANOVA\par}
  \vspace{0.7cm}
  {\color{TecRojo}\rule{\textwidth}{1.2pt}}
  \vspace{0.8cm}

  Facultad MA1001B\\[0.3cm]
  Periodo\\[0.3cm]
  Tecnol\'ogico de Monterrey
\end{frame}

\begin{frame}{La pregunta de supuestos}
  \begin{alertblock}{Pregunta guía}
    Tras un $F$ significativo y un seguimiento Tukey, ¿qué debe seguir
    siendo cierto de residuos, dispersiones y el diseño de muestreo?
  \end{alertblock}

  \vspace{0.6em}
  Al final de esta lección, usted puede:
  \begin{itemize}
    \item formar residuos $e_{ij}=y_{ij}-\bar y_i$ e inspeccionar su
      histograma;
    \item usar la prueba de Levene y gráficas residuo-frente-a-ajustado;
    \item tratar la independencia como una afirmación de diseño;
    \item explicar por qué $n$ igual no rescata una heterocedasticidad
      fuerte.
  \end{itemize}

  \vfill
  \scriptsize Todos los tiempos de ciclo usados hoy son sintéticos. No se usan datos reales de empresa.
\end{frame}

\begin{frame}{Tres supuestos, un residuo}
  \small
  \begin{center}
    \begin{tabular}{ll}
      \toprule
      \textbf{Supuesto} & \textbf{Qué chequeamos} \\
      \midrule
      Normalidad & Histograma de $e_{ij}=y_{ij}-\bar y_i$ \\
      Varianzas iguales & Levene y residuo frente a ajustado \\
      Independencia & Estaciones aleatorizadas, no una gráfica \\
      \bottomrule
    \end{tabular}
  \end{center}

  \vspace{0.5em}
  \begin{exampleblock}{Experimento de empaque equilibrado}
    $n=12$ estaciones en cada método. La media residual es $0.000000$ por
    construcción dentro de cada grupo.
  \end{exampleblock}
\end{frame}

\begin{frame}[fragile]{Paso 1: Normalidad de residuos}
  \begin{columns}[T]
    \begin{column}{0.40\textwidth}
      \small
      \begin{lstlisting}
e = y - group_mean
W, p = shapiro(e)
      \end{lstlisting}

      \begin{block}{Salida verificada}
        DE residual $2.630914$\\
        Shapiro $W=0.977891$\\
        $p=0.673807$
      \end{block}
    \end{column}
    \begin{column}{0.60\textwidth}
      \centering
      \includegraphics[width=\textwidth]{../figuras/diagnostico_anova_01_normalidad_residuos.png}
      \vspace{0.2em}
      \scriptsize Histograma de residuos del Paso 1
    \end{column}
  \end{columns}
\end{frame}

\begin{frame}{La dispersión igual es un segundo chequeo, separado}
  \small
  DE muestrales: $3.067311$ (Estándar), $2.353358$ (Guiado), $2.660249$
  (Automatizado).
  \[
    \text{Levene }W=0.457821,\qquad p=0.636619.
  \]
  Residuo frente a ajustado no debería abrirse en abanico. La
  independencia se hereda de la asignación aleatoria de estaciones: una
  gráfica de residuos no puede probarla.
\end{frame}

\begin{frame}[fragile]{Paso 2: Levene y residuo frente a ajustado}
  \begin{columns}[T]
    \begin{column}{0.40\textwidth}
      \small
      \begin{lstlisting}
W, p = levene(s, g, a)
e = y - fitted
      \end{lstlisting}

      \begin{block}{Salida verificada}
        Levene $p=0.636619$\\
        DE $3.067$, $2.353$, $2.660$\\
        Sin evidencia de dispersión desigual
      \end{block}
    \end{column}
    \begin{column}{0.60\textwidth}
      \centering
      \includegraphics[width=\textwidth]{../figuras/diagnostico_anova_02_levene_homocedasticidad.png}
      \vspace{0.2em}
      \scriptsize Gráfica residuo-frente-a-ajustado del Paso 2
    \end{column}
  \end{columns}
\end{frame}

\begin{frame}{Paso 3: $n$ igual no rescata una heterocedasticidad fuerte}
  \begin{columns}[T]
    \begin{column}{0.42\textwidth}
      \small
      Sigue $n=12$ por método. DE $1.534$, $1.324$, $9.976$.

      \vspace{0.4em}
      Levene $p=2.455155\times10^{-4}$. ANOVA $F=9.678598$,
      $p=4.924585\times10^{-4}$.

      \vspace{0.5em}
      \textbf{Límite:} ANOVA de $n$ igual no es inmune a una desigualdad
      fuerte de varianzas. La Lección 47 pasa de medias de grupo a una
      recta ajustada.
    \end{column}
    \begin{column}{0.58\textwidth}
      \centering
      \includegraphics[width=\textwidth]{../figuras/diagnostico_anova_03_limite_varianzas_desiguales.png}
      \vspace{0.2em}
      \scriptsize Ejemplo heterocedástico de $n$ igual del Paso 3
    \end{column}
  \end{columns}
\end{frame}

\begin{frame}{Verificación de conceptos}
  \begin{enumerate}
    \item ¿Por qué la media residual es exactamente 0 dentro de cada método
      de empaque?
    \item ¿Qué dice Levene $p=0.636619$, y qué no dice?
    \item ¿Por qué una gráfica de residuos no puede probar independencia?
    \item En el Paso 3, $n=12$ en cada grupo. ¿Por qué el $F$ significativo
      sigue siendo poco confiable?
  \end{enumerate}

  \vspace{0.6em}
  \begin{block}{Conexión con la Lección 47}
    La sesión puede continuar directamente con regresión lineal simple,
    mínimos cuadrados y $R^2$.
  \end{block}

  \vfill
  \scriptsize Fuente: OpenStax, \textit{Introductory Business Statistics 2e},
  Secciones 12.2 y 12.4. CC BY-NC-SA.
\end{frame}

\appendix



\end{document}
